\documentstyle[%


righttag%


,11pt%


,amssymb%


,amscd%


,russian%               remove in Eng.


,izvru%                 Set izven% in Eng.


]{amsart}





%


\newcommand{\Id}{\operatorname{Id}}


\newcommand{\ch}{\operatorname{ch}}


\newcommand{\sh}{\operatorname{sh}}


\def\pd{{\partial}}


\newcommand{\Pf}{\operatorname{Pf}}


\newcommand{\Ann}{\operatorname{Ann}}


\def\R{{\Bbb R}}


\def\C{{\Bbb C}}


\def\SS{{\Bbb S}}


\def\T{{\Bbb T}}


\def\Z{{\Bbb Z}}


\def\K{{\Bbb K}}


\def\rphi{\varphi}


\def\reps{\varepsilon}


\def\cA{{\cal A}}


\def\cD{{\cal D}}


\def\cF{\cal F}


\def\cG{{\cal G}}


\def\cI{{\cal I}}


\def\cL{{\cal L}}


\def\cM{{\cal M}}


\def\cO{{\cal O}}


\def\cP{{\cal P}}


\def\cS{{\cal S}}


\def\cT{{\cal T}}


\def\cU{{\cal U}}


\def\g{{\frak g}}


\def\gl{{\frak  gl}}


\def\gK{{\frak  K}}


\def\sp{{\frak  sp}}


\def\gU{{\frak  U}}


\def\gV{{\frak  V}}


\def\gX{{\frak  X}}


\def\sp{{\frak  sp}}


\def\rest#1#2{#1 |_{#2}}


\def\germ#1#2{\langle #1 \rangle_{#2}}


\def\fref#1{(\ref{#1})}


\newcommand{\tts}[1]{}





\theoremstyle{plain}


\theorembodyfont{\it}


\newtheorem{statnonum}{\hskip\parindent Утверждение} %% !!


\renewcommand{\thestatnonum}{}





\newtheorem{lemnonum}{\hskip\parindent Лемма}


\renewcommand{\thelemnonum}{}





\theoremstyle{definition}


\theorembodyfont{\rm}


\newtheorem{examnonum}{\hskip\parindent Пример}


\renewcommand{\theexamnonum}{}





%\hoffset=-15mm%                  For Eng.


\setcounter{page}{45}


\god{2004}


\nomer{11~(510)} %                        Remove in Eng.


%\nomer{Vol.\,48, No.\,11}%               Set in Eng.


%\str{\pageref{begin}--\pageref{end}}%  Set in Eng.


\udk{514.763}





\author{\it М.А.\,Малахальцев}


% NB:   ^^ remove in Eng.


\title{Пучок локальных гамильтонианов\\  на симплектическом


многообразии\\  с особенностями Мартине}





\date{}


%\pagestyle{myheadings}%         Set in Eng.


\pagestyle{plain} %              Remove in Eng.


\begin{document}


%\thispagestyle{plain}%          Set in Eng.


\maketitle


\label{begin}








\section{Введение}





Пусть $\omega_0$ есть замкнутая 2-форма на $\R^{2n}$. Симплектическое


многообразие с особенностями типа $\omega_0$ есть $2n$-мерное


многообразие $M$, наделенное замкнутой формой $\omega$ такой, что


для любой точки $p \in M$ существует окрестность


$U(p)$ и диффеоморфизм $\phi_p : U(p) \to \R^{2n}$ со свойством


$\rest{\omega}{U(p)} = \phi^*_p\omega_0$ (см. \cite{MA}).





В \cite{Mart1}--\cite{GolTish}


доказано, что на $\R^4$ есть пять типов ростков общего положения замкнутой


2-формы. Из них четыре являются структурно устойчивыми ~---


это ростки следующих форм в начале координат:


\smallskip





тип 0


$$


\omega_0 = dx^1\wedge dx^2 + dx^3\wedge dx^4,


$$





тип I


$$


\omega_0 = x^1 dx^1\wedge dx^2 + dx^3\wedge dx^4,


$$





тип II-e (эллиптический тип)


\begin{equation}


\omega_0 = dx^1 \wedge dx^2 + x^3 dx^1 \wedge dx^4 + x^3 dx^2 \wedge dx^3


+ x^4 dx^2 \wedge dx^4 + (x^1- (x^3)^2) dx^3 \wedge dx^4,


\label{eq3}


\end{equation}





тип II-h (гиперболический тип)


\begin{equation}


\omega_0 = dx^1 \wedge dx^2 + x^3 dx^1 \wedge dx^4 + x^3 dx^2 \wedge dx^3


- x^4 dx^2 \wedge dx^4 + (x^1- (x^3)^2) dx^3 \wedge dx^4.


\label{eq4}


\end{equation}





Данные типы можно охарактеризовать следующим образом (подробности см. в


\cite{Mart1}). Пусть $\omega$ есть замкнутая 2-форма на четырехмерном


многообразии $M$. Предположим, что множество точек, в которых форма


$\omega$


вырождена, образует трехмерное подмногообразие $\Sigma$ в $M$.





Если точка $p$ не лежит в $\Sigma$, то по теореме Дарбу в


окрестности точки $p$ \ 2-форма $\omega$ приводится к виду 0.





Пусть точка $p$ лежит в $\Sigma$ и $\omega(p) \ne 0$. Если аннулятор


$E_p$


формы $\omega(p)$ трансверсален $\Sigma$, то форма $\omega$


приводится в некоторой окрестности $U$ точки $p$ к типу I. При этом


открытое в $\Sigma$ множество $V = U \cap \Sigma$ состоит из точек


типа I, а всюду плотное в $U$ подмножество $U \setminus V$ состоит из


точек типа 0.





Пусть $\Sigma' \subset \Sigma$~--- множество точек $p$ таких, что


аннулятор $E_p \subset T_p \Sigma$. Если $\Sigma'$ в окрестности точки


$p$ есть


подмногообразие и $E$ трансверсально $\Sigma'$,


то форма $\omega$ приводится в некоторой окрестности $U$ точки $p$


к типу II-e или II-h. Пусть $V = U \cap \Sigma$,


$W = U \cap \Sigma'$.


Тогда всюду плотное в $U$ множество $U \setminus \Sigma$ состоит из


точек типа 0, всюду плотное в $V$ множество $V \setminus W$ состоит из


точек типа I, и множество $W$ состоит из точек типа II.





Пусть на четырехмерном многообразии задана замкнутая 2-форма


$\omega$. Тогда если все точки $M$ имеют тип 0, то $(M,\omega)$~---


симплектическое многообразие. Если любая точкa $M$ имеет тип 0 или I, то


скажем, что $(M,\omega)$~--- {\em симплектическое многообразие с


особенностями Мартине $\mathrm I$ типа}.


Если любая точкa $M$ имеет тип 0, I, или II, то


скажем, что $(M,\omega)$~--- {\em симплектическое многообразие с


особенностями Мартине $\mathrm{II}$ типа}.





Заметим, что симплектическое многообразие $(M,\omega)$ с особенностями


Мартине есть


симплектическое многообразие с особенностями типа $\omega_0$,


где $\omega_0$ задана формулами (\ref{eq3}) или


(\ref{eq4}).


Действительно, пусть $\phi(p) = (x^1_0, x^2_0, x^3_0, x^4_0)$.


Тогда точка $p$ имеет тип 0 тогда и только тогда, когда


$x^1_0\ne 0$; тип I тогда и только тогда, когда


$x^1_0=0$ и $(x^3_0)^2 + (x^4_0)^2 \ne 0$; тип II тогда и только


тогда, когда $x^1_0=0$, $x^3_0=0$ и $x^4_0=0$. Таким образом, симплектическая


структура есть частный случай симплектической структуры с


особенностями Мартине типа I, которая, в свою очередь, есть частный


случай симплектической структуры с особенностями Мартине типа II.








Симплектические многообразия $(M,\omega)$ с особенностями Мартине типа I


(под названием ``folded symplectic structures'') изучались


в работе \cite{SGW}, где найдены условия, при которых $\omega$ есть


прообраз симплектической формы на симплектическом орбифолде и


показано, что $\omega$ индуцирует единственную с точностью до


гомотопии спинорную структуру. Отметим, что в этой работе


рассматриваются многообразия произвольной четной размерности, а форма


$\omega_0$ имеет вид


\begin{equation}


 \label{eq5}


x^1 dx^1 \wedge dx^2 + dx^3 \wedge dx^4 + \dotsb + dx^{2n-1}\wedge dx^{2n}.


\end{equation}





В \cite{S} показано, что на


любом компактном ориентируемом четырехмерном многообразии существует


симплектическая структура с особенностями Мартине типа I. В \cite{D1}


доказаны аналоги теоремы о несжимаемости объема для $\R^{2n}$ с формой


(\ref{eq5}).





Пусть $(M,\omega)$ есть симплектическое многообразие с особенностями


типа $\omega_0$. Векторное поле $V$ на открытом подмножестве $U$


называется {\em гамильтоновым векторным полем}, если $L_V \omega =0$, где


$L_V$ есть производная Ли в направлении $V$. Пусть $\gX_h(U)$ есть


алгебра Ли гамильтоновых векторных полей на $U$. Ясно, что $U \to


\gX(U)$ есть пучок алгебр Ли на $M$. В данной работе изучается пучок


гамильтоновых векторных полей на симплектическом многообразии с


особенностями Мартине.





Все рассматриваемые в работе многообразия и отображения предполагаются


гладкими. Используются следующие обозначения: $C^\infty_M$~--- пучок


гладких функций на многообразии $M$, $\gX_M$~--- пучок векторных


полей, $\Omega^*_M$~--- пучок дифференциальных форм. Для пучка $\cG$


на многообразии $M$ через $\cG_p$ будем обозначать множество ростков


этого пучка в точке $p \in M$.








\section{Основные результаты}





Пусть $(M,\omega)$ есть симплектическое многообразие с особенностями


типа $\omega_0$. Гладкая функция $f : U \to \R$, где $U \subset M$~---


открытое множество, называется {\em локальным гамильтонианом}


(ср. \cite{Fom}),


если существует векторное поле $V$ на $U$ такое,


что $i_V\omega = df$. Из свойства производной Ли $L_V = d i_V + i_V d$


следует $V \in \gX_h(U)$. Ясно, что локальные гамильтонианы


образуют пучок, обозначим его $\cT$.





В \cite{MA} было показано, что имеет место точная последовательность


пучков


$$


0 \to \R_M @>i>> \cT @>{\pi}>> \gX_\Gamma \to 0,


$$


где $\R_M$ есть пучок локально постоянных функций на $M$, $i$~---


включение


$\R_M$ в $\cT$, и $\pi(f) = V$, если $df = i_V \omega$, т.\,е. $\pi$


отображает гамильтониан в соответствующее гамильтоново векторное поле.


Отсюда следует, что существует точная последовательность когомологий


многообразия $M$ с коэффициентами в пучках $\R_M$, $\cT$ и $\gX_h$:


\begin{equation} \label{eq7}


\dotsb \to


\check H^0(M;\gX_\Gamma) @>{\delta}>>


\check H^1(M;\R_M) @>{i_*}>>


\check H^1(M;\cT) @>{\pi_*}>>


%%\\


%%\overset{\pi_*}{\to}


\check H^1(M;\gX_\Gamma) @>{\delta}>>


\check H^2(M;\R_M) @>{i_*}>> \dotsb


\end{equation}





В дальнейшем под симплектической структурой с особенностями Мартине


мы будем понимать симплектическую структуру с особенностями типа


$\omega_0$, где $\omega_0$ имеет вид (\ref{eq3}) или (\ref{eq4}),


на четырехмерном многообразии.


В этом параграфе найдем когомологии пучка локальных гамильтонианов


на симплектическом многообразии с особенностями Мартине.








Пусть $\mu$~--- форма объема на $M$. Тогда


$\omega \wedge \omega = \Pf_\mu(\omega) \mu$, где $\Pf_\mu(\omega)$


называется


{\em пфаффианом формы $\omega$ относительно $\mu$}.


Очевидно, если $\mu' = \lambda \mu$~--- другая форма объема, где


$\lambda$~--- функция, не обращающаяся в нуль, то


$\Pf_\mu(\omega) = \lambda \Pf_{\mu'}(\omega)$. Следовательно, в пучке


$C^\infty$ колец гладких функций определен {\em подпучок $\cI$ главных


идеалов, порожденных $\Pf_\mu(\omega)$}:


$$


 \cI(U) = \{\rest{\Pf_\mu(\omega)}{U} \cdot f \mid f \in C^\infty(U)\},


$$


который не зависит от выбора формы объема $\mu$.





Далее заметим, что {\em $\cI$ есть в точности пучок идеалов, сечениями


которого являются функции, обращающиеся в нуль на $\Sigma$, т.\,е.


$C^\infty/\cI$ есть пучок ростков гладких функций на $\Sigma$}.








\begin{lemnonum} \label{lm1}


Пусть $(M,\omega)$ есть симплектическое многообразие с особенностями


Мартине. Пусть ковекторное поле $\xi$ определено на открытом множестве


$U \subset M$. Тогда $\xi = i_V \omega$ для некоторого гладкого векторного


поля


$V$ на $U$ тогда и только тогда, когда для любой точки $p \in U$


ковектор $\xi(p)$ обращается в нуль на аннуляторе $\omega(p)$. При


этом векторное поле $V$ определяется единственным образом.





В частности, $f$ есть локальный гамильтониан тогда и только тогда,


когда $df$ обращается в нуль на аннуляторе $\omega$.


\end{lemnonum}








\begin{pf}


Вначале положим $U = M$.


Пусть $\alpha_\omega : TM \to T^*M$, $\alpha(V) = i_V\omega$, --- морфизм


векторных расслоений, определяемый 2-формой $\omega$. Нам надо


доказать, что сечение $\xi$ расслоения $T^*M$, определенное на $U$,


есть образ сечения


$V$ расслоения $TM$ при отображении $\alpha_\omega$ тогда и только тогда,


когда $\xi$


обращается в нуль на аннуляторе $\omega$.





Если $\xi = \alpha_\omega(V)$, то из определения аннулятора немедленно


следует, что $\xi$ обращается на нем в нуль.





Докажем обратное утверждение.


Прежде всего построим кососимметричное тензорное поле $\wt\omega \in


\Omega_2(M)$ такое, что


\begin{equation}


 \label{eq12}


 \wt\omega{}^{sk}\omega_{sm} = \Pf_\mu(\omega)\delta^k_m,


\end{equation}


где $\mu$~--- некоторая форма объема.


Зафиксируем некоторую метрику $g$


на $M$ и пусть $\mu \in \Omega^4(M)$~--- соответствующая форма


объема. Тогда, как известно, существует тензорное поле $\wt\mu \in


\Omega_4(M)$


такое, что в локальных координатах $\wt\mu{}^{ijkl} \mu_{ijkm} =


\delta^l_m$. Положим $\wt\omega{}^{kl} = \wt\mu{}^{ijkl}\omega_{ij}$.


Используя определение $\Pf_\mu(\omega)$, а также кососимметричность


формы $\omega$


и тензора $\wt\mu$, запишем следующие равенства:


\begin{multline*}


\Pf_\mu(\omega)\delta^p_m = \Pf_\mu(\omega) \wt\mu{}^{ijkp} \mu_{ijkm} =


\wt\mu{}^{ijkp} \omega_{[ij}\omega_{km]} =


\\


=\frac{1}{3}(\wt\mu{}^{ijkp}\omega_{ij}\omega_{km} +


\wt\mu{}^{ijkp}\omega_{ik}\omega_{mj} +


\wt\mu{}^{ijkp}\omega_{im}\omega_{jk}) =


\wt\mu{}^{ijkp}\omega_{ij}\omega_{km} = \wt\omega{}^{kp}\omega_{km}.


\end{multline*}


Таким образом, построенное тензорное поле $\wt\omega \in \Omega_2(M)$


обладает требуемым свойством (\ref{eq12}). Тензорное поле $\wt\omega$


определяет морфизм векторных расслоений


$\wt\alpha_\omega : T^*M \to TM$, причем $\wt\alpha_\omega \alpha_\omega =


\Pf_\mu(\omega) \Id_{TM}$.





Пусть $\xi \in \cI\cdot\Omega^1(M)$, т.\,е. $\xi =


\Pf_\mu(\omega)\eta$, где $\eta \in \Omega^1(M)$.


Положим $V = \wt\alpha_\omega \eta \in \gX(M)$. Тогда


$$


\alpha_\omega(V) = \alpha_\omega \wt\alpha_\omega(\eta) =


\Pf_\mu(\omega)\eta = \xi.


$$


Таким образом, если сечение $\xi$ векторного расслоения $T^*M$


обращается в нуль на $\Sigma$, то существует векторное поле $V$ на


$M$ такое, что $\alpha_\omega(V) = \xi$.








Пусть теперь $\xi$ есть некоторое ковекторное поле на $M$,


обращающееся в нуль на аннуляторе $\omega$.


Так как $\omega$ есть симплектическая форма с особенностями


Мартине, морфизм векторных расслоений


$\alpha_\omega : \rest{TM}{\Sigma} \to \rest{T^*M}{\Sigma}$ имеет


постоянный ранг 2, и его образ состоит из сечений расслоения


$\rest{T^*M}{\Sigma}$, обращающихся в нуль на аннуляторе $\omega$.


Поэтому $\rest{\xi}{\Sigma} = \alpha_\omega(\wt V)$, где $\wt V$ есть


сечение расслоения $\rest{TM}{\Sigma}$. Теперь продолжим


$\wt V$ до векторного поля $V$ на $M$ ($\Sigma$ замкнуто


в $M$), тогда $\xi' = \xi - \alpha_\omega(V)$ обращается в нуль на


$\Sigma$.


По доказанному ранее $\xi' = \alpha_\omega(V'')$ для некоторого


векторного поля $V''$ на $M$, поэтому


$$


\xi = \xi' + \alpha_\omega(V'') =


\alpha_\omega(V)+\alpha_\omega(V'')=\alpha_\omega(V+V'').


$$





Таким образом, мы показали, что для любого $\xi \in \Omega^1(M)$,


обращающегося в нуль на аннуляторе формы $\omega$, существует $V \in


\gX(M)$ такое, что


$\xi = \alpha_\omega(V)$. Так как


$\alpha_\omega$ есть изоморфизм на всюду плотном множестве точек,


получаем, что векторное поле $V$ определено однозначно.





Осталось заметить, что открытое подмногообразие симплектического


многообразия с особенностями Мартине само есть симплектическое


многообразие с


особенностями Мартине. Таким образом, лемма доказана для открытого


подмножества $U \subset M$.


\end{pf}








На подмногообразии $\Sigma$ определено слоение $\cF$ с особенностями,


регулярные слои которого суть интегральные кривые одномерного


распределения на $\Sigma \setminus \Sigma'$, полученного пересечением


аннулятора $\omega$ с касательными пространствами к $\Sigma$, а особым


слоем является кривая $\Sigma'$. Обозначим через $\cF_b$ пучок базовых


функций слоения $\cF$, т.\,е. пучок функций, локально


постоянных вдоль слоев $\cF$.





Пусть $X$~--- топологическое пространство, $A \subset


X$~--- замкнутое подпространство, и $\cG$~--- пучок колец на


$A$. Тогда определен пучок $\cG^X$ на $X$, порожденный предпучком:


$U \to {0}$, если $U \cap A = \emptyset$, и $U \to \cG(U \cap A)$ в


противном случае (см. \cite{Bre}).








Пусть $r: C^\infty_M \to (C^\infty_\Sigma)^M$ есть


морфизм пучков, порожденный ограничением функций на $\Sigma$, т.\,е.


если $U \cap \Sigma = \emptyset$, то $(C^\infty_\Sigma)^M(U) = 0$ и


$r_U : C^\infty_M(U) \to (C^\infty_\Sigma)^M(U)$ есть нулевой гомоморфизм.


Если $U \cap \Sigma = U' \ne \emptyset$, то


$r_U : C^\infty_M(U) \to C^\infty_\Sigma(U')$ есть ограничение функции


$f \in C^\infty(U)$ на $U'$.


По определению $\cI^2$ есть подпучок пучка $C^\infty_M$,


включение $\cI^2 \to C^\infty_M$ обозначим через $i$,


также положим $\wt\cF_b = \cF_b^M$.





\begin{statnonum} \label{st1}


Имеет место точная последовательность пучков на $M$


\begin{equation} \label{eq9}


0 \to \cI^2 @>{i}>> \cT @>{r}>> \wt\cF_b \to 0.


\end{equation}


\end{statnonum}





\begin{pf}


Прежде всего заметим, что если $f \in \cI^2$, то $df = 0$ в точках


$\Sigma$, поскольку~$\cI$ есть пучок идеалов функций обращающихся в нуль


на $\Sigma$.


Следовательно, по лемме $df = i_V\omega$, значит, $f$ есть


сечение пучка $\cT$. Таким образом, $i : \cI^2 \to C^\infty_M$ отображает


$\cI^2$ в $\cT$.





Покажем, что морфизм пучков


$r : C^\infty_M \to (C^\infty_\Sigma)^M$ отображает подпучок $\cT$ пучка


$C^\infty_M$ в


подпучок $\wt\cF_b$ пучка $(C^\infty_\Sigma)^M$.


Для этого достаточно показать, что для любой $p \in M$ морфизм пучков $r$


отображает кольцо ростков $\cT_p$ в кольцо ростков $(\wt\cF_b)_p$.





Если $p \notin \Sigma$, то требуемое утверждение непосредственно


следует из определений. Пусть теперь $p$ лежит в $\Sigma$ и $\germ{f}{p}$


лежит в $\cT_p$. Так как $df$ обращается в нуль на аннуляторе $E$ формы


$\omega$, $df$ равно нулю на $T\Sigma \cap E$, следовательно, функция $f$


постояннa вдоль любого неособого слоя слоения $\cF$ на


$\Sigma$. Таким образом, если $p \not\in \Sigma'$, то $\germ{f}{p}$


есть росток базовой функции слоения $\cF$, т.\,е. $\germ{f}{p} \in


\wt\cF_b$.





Пусть $p \in \Sigma'$. Тогда в окрестности точки $p$ можно взять


систему координат, в которой форма $\omega$ имеет вид (\ref{eq3}) или


(\ref{eq4}), а точка $p$ имеет нулевые координаты.


Рассмотрим эллиптический случай (\ref{eq3}).


В этой системе координат $\Sigma$ задается уравнением $x^1=0$,


а $\Sigma'$ уравнением $x^1=x^3=x^4=0$. Далее, аннулятор $E$ натянут


на векторы


\begin{equation}


\label{eq11}


V_1 = x^3\pd_1 + \pd_3, \quad V_2 = x^4\pd_1 - x^3\pd_2 +\pd_4,


\end{equation}


а слоение $\cF$ порождено векторным полем


$$


W = -(x^3)^2\pd_2 - x^4\pd_3 + x^3\pd_4.


$$


Из условия $\rest{df}{E}=0$ следует, что при $x^1=0$


\begin{equation}


\label{eq14}


\begin{split}


x^3 \pd_1 f(0,x^2,x^3,x^4) + \pd_3 f(0,x^2,x^3,x^4) &= 0,


\\


x^4 \pd_1 f(0,x^2,x^3,x^4) - x^3 \pd_2 f(0,x^2,x^3,x^4) + \pd_4


f(0,x^2,x^3,x^4) &= 0.


\end{split}


\end{equation}


Дифференцируя первое уравнение из \eqref{eq14} по $x^4$, а второе по


$x^3$ и используя равенство\\ $\pd_{34}f(0,x^2,x^3,x^4) =


\pd_{43}f(0,x^2,x^3,x^4)$, получаем


$$


x^3 \pd_{14} f(0,x^2,x^3,x^4) = x^4 \pd_{13} f(0,x^2,x^3,x^4) -


\pd_2 f(0,x^2,x^3,x^4) - x^3 \pd_{23} f(0,x^2,x^3,x^4).


$$


Подставляя в последнее соотношение $x^3=0$, $x^4=0$, получим


$\pd_2 f(0,x^2,0,0) = 0$, следовательно, функция $f$ постоянна вдоль особого


слоя. Таким образом, $\germ{f}{p} \in \wt\cF_b$.





Гиперболический случай разбирается аналогично. Подмногообразия


$\Sigma$ и $\Sigma'$ задаются теми же уравнениями, а аннулятор $E$


натянут на векторные поля


\begin{equation}


\label{eq13}


V_1 = x^3\pd_1 + \pd_3, \quad V_2 = - x^4\pd_1 - x^3\pd_2 + \pd_4.


\end{equation}


Имеем систему, аналогичную (\ref{eq14}), и такими же рассуждениями


получаем, что $f$ постоянна вдоль особого слоя. Таким образом, и в этом


случае $\germ{f}{p} \in \wt\cF_b$.








Перейдем к доказательству точности последовательности (\ref{eq9}).


Нам надо доказать, что для любой точки $p \in M$ последовательность


\begin{equation} \label{eq10}


0 \to (\cI^2)_p @>{i_p}>> \cT_p @>{r_p}>> (\wt\cF_b)_p \to


0


\end{equation}


точна. Прежде всего заметим, что для любой $p \not\in \Sigma$ эта


последовательность точна, т.\,к. $(\cI^2)_p = (C^\infty_M)_p$,


$\cT_p = (C^\infty_M)_p$, $(\wt\cF_b)_p=0$, и $i_p$ есть тождественное


отображение.





Далее, для любой точки $p$ точность последовательности (\ref{eq10}) в


члене $(\cI^2)_p$ следует из определений. Следовательно, надо доказать


ее точность в оставшихся двух членах. Зафиксируем некоторую форму объема


и для краткости


обозначим $\Pf_\mu(\omega)$ через $\lambda$. Заметим, что $d\lambda


\ne 0$ в точках $\Sigma$.





{\em Точность в члене $\cT_p$.}


Из определений немедленно следует, что $r_p i_p = 0$. Теперь, пусть


$\germ{f}{p} \in \cT_p$ и $r_p(\germ{f}{p})=0$. Тогда


$\germ{f}{p} \in \cI_p$, следовательно, в некоторой окрестности $U$


точки $p$ имеет место $f = \lambda g$, где $g \in C^\infty(U)$.


Тогда на $U' = U \cap \Sigma$ имеем $df = g d\lambda$. Так как $f \in


\cT(U)$, то


$\rest{df}{E} = 0$. Но в плотном подмножестве точек


$U'' = U' \setminus (U \cap \Sigma')$ множества $U'$ аннулятор $E$


трансверсален


$T\Sigma$, следовательно, $\rest{g}{U''} = 0$, т.\,к. $d\lambda \ne


0$. Значит, $\rest{g}{U'} = 0$, следовательно, $g = h \lambda$, откуда


$f = g \lambda = h \lambda^2$. Таким образом, $\germ{f}{p} \in (\cT^2)_p$.





{\em Точность в члене $(\wt\cF_b)_p$.}


Покажем, что для любой точки $p$ в $\Sigma$ и


любого $\germ{f}{p} \in (\wt\cF_b)_p$ существует


$\germ{\wh f\,}{p} \in \cT_p$ такой, что $r_p(\germ{\wh f\,}{p}) =


\germ{r(\wh f\,)}{p} = \germ{f}{p}$.





Пусть $p \in \Sigma\setminus\Sigma'$. Тогда можно выбрать


представителя $f$ ростка $\germ{f}{p} \in (\wt\cF_b)_p$, определенного


нa открытом множестве $U' = U \cap \Sigma$, не пересекающем $\Sigma'$.


Следовательно, на $U'$ аннулятор $E$ формы $\omega$ трансверсален $T\Sigma$, и


пусть $l = E \cap T\Sigma$~--- одномерное распределение,


задающее ограничение слоения $\cF$ на $U'$.


Выберем поле репера $W_1, W_2$ расслоения $E$ над $U$ так, чтобы


$W_1(q) \in l(q)$ и $d\lambda(W_2(q))=1$ для любого $q \in U'$.





Пусть $\wt f$ есть некоторое продолжение $f$ на $U$, и пусть


$g = d\wt f(W_2) : U' \to \R$. Обозначим через $\wt g$ некоторое


продолжение $g$ на $U$ и рассмотрим функцию


$\wh f = \wt f - \wt g \lambda$ на $U$.


Так как $\lambda(q)=0$ для любой $q \in U'$, то $r(\wh f) = f$.


Далее, в точках $U'$ имеем $d\wh f = d f + g d\lambda$. Cледовательно,


$d\wh f(W_1) = 0$, т.\,к. $f$~--- базовая функция слоения $\cF$ и


$d\lambda(W_1)= 0$, и $d\wh f(W_2) = d\wt f(W_2) - g d\lambda(W_2) =


0$ в силу выбора $W_2$ и определения функции $g$. Таким образом,


$d\wh f$ обращается в нуль на $E$, следовательно, $\wh f \in \cT(U)$.





Теперь докажем точность в точках $p \in \Sigma'$. Вначале рассмотрим


эллиптический случай. Найдем вид


базовых функций слоения $\cF$ в канонической системе координат


(\ref{eq3}), заданной в окрестности $U'$ точки $p$.


Перейдем к цилиндрической системе координат на $U'\setminus(U'\cap


\Sigma')$:


$x^2=u$, $x^3= \rho \cos


\varphi$, $x^4 = \rho \sin \varphi$. Тогда распределение $l$ натянуто


на векторное поле $W= \pd_\varphi - \rho^2 \cos^2\varphi\, \pd_2$ и из


условия $Wf=0$ получаем, что при $\rho >0$


$$


f(x^2,\rho\cos\varphi,\rho\sin\varphi) = F(x^2 + \tfrac{1}{2}\rho^2(\varphi


+\tfrac{1}{2}\sin 2\varphi)) + G(\rho),


$$


где $F(p)$~--- гладкая функция, $G(\rho)$~--- гладкая функция,


определенная при $\rho>0$. Подставив в это равенство $\varphi=0$,


получим


$$


f(x^2,\rho,0) = F(x^2) + G(\rho).


$$


Отсюда $\pd_2 f(x^2,\rho,0) = F'(x^2)$. Это равенство имеет место при


$\rho>0$, но $\pd_2 f(x^2,\rho,0)\to 0$ при $\rho \to 0$, т.\,к.


$\pd_2 f(x^2,0,0)=0$ по определению базовой функции.


Таким образом, $F(x^2)$ постоянна и


$$


f(x^2,\rho\cos\varphi,\rho\sin\varphi) = G(\rho)+C,


$$


где $C$~--- постоянная. В силу гладкости $f(x^2,x^3,x^4)$ отсюда


следует $f(x^2,x^3,x^4) = g((x^3)^2 + (x^4)^2) + C$ для некоторой


гладкой функции $g(u)$. Положим


$$


\wh f(x^1,x^2,x^3,x^4) = C + g((x^3)^2 + (x^4)^2) - 2 x^1 g'((x^3)^2 +


(x^4)^2).


$$


Тогда при $x^1=0$ имеем $d\wh f = 2 g'((x^3)^2 + (x^4)^2) \theta$,


где $\theta = dx^1 + x^3 dx^3 + x^4 dx^4$. Поэтому $d\wh f$


обращается в нуль на векторных полях (\ref{eq11}), задающих базис $E$,


следовательно, по лемме $\wh f$ лежит в $\cT(U)$.


Наконец, ясно, что $r(\wh f)=f$, что доказывает сюръективность $r_p$.


Таким образом, в эллиптическом случае точность в члене $(\wt\cF_b)_p$ для


точек $p \in


\Sigma'$ доказана.





Гиперболический случай рассматривается аналогично. Вначале находим вид


базовых функций слоения $\cF$ в канонической системе координат


(\ref{eq4}), заданной в окрестности $U'$ точки $p$, для чего переходим к


гиперболической цилиндрической системе координат на


$U'\setminus(U'\cap \Sigma')$:


$x^2=u$, $x^3= \rho \ch \varphi$, $x^4 = \rho \sh \varphi$.


Получаем, что для базовой функции $f$ при $\rho >0$ имеет место равенство


$$


f(x^2,\rho\cos\varphi,\rho\sin\varphi) =


F(x^2 + \tfrac{1}{2}\rho^2(\varphi + \tfrac{1}{2}\sh 2\varphi)) + G(\rho),


$$


где $F(p)$~--- гладкая функция, $G(\rho)$~--- гладкая функция,


определенная при $\rho>0$. Далее, рассуждая аналогично эллиптическому


случаю, получаем $f(x^2,x^3,x^4) = g((x^3)^2 - (x^4)^2) +


C$. Положим


$$


\wh f(x^1,x^2,x^3,x^4) = C + g((x^3)^2 - (x^4)^2) - 2 x^1 g'((x^3)^2 -


(x^4)^2).


$$


Ясно, что $d\wh f$ обращается в нуль на векторных полях (\ref{eq13}),


задающих базис $E$,


следовательно, по лемме $\wh f$ лежит в $\cT(U)$. Кроме


того, $r(\wh f)=f$, таким образом, $r_p$ сюръективно.


Итак, и в гиперболическом случае точность в члене $(\wt\cF_b)_p$ для


точек $p \in


\Sigma'$ доказана.


\end{pf}





\begin{cor} \label{cor2}


Для симплектического многообразия с особенностями Мартине


имеет место изоморфизм групп когомологий


$$


H^q(M;\cT) \cong H^q(\Sigma;\cF_b), \quad q> 0,


$$


где $\cT$~--- пучок локальных гамильтонианов на $M$, $\cF_b$~--- пучок


базовых функций слоения с особенностями, индуцированного на особом


подмногообразии $\Sigma$ аннулятором формы $\omega$.


\end{cor}





\begin{pf}


Заметим, что пучок $\cI^2$ тонкий, имеем $H^q(M,\cI^2)\cong 0$, $q > 0$.


В силу предложения \ref{st1} имеет место точная последовательность


$$


\dotsb \to


\check H^q(M;\cI^2) @>{i}>>


\check H^q(M;\cT) @>{r}>>


\check H^q(M;\wt F_b) \to


\check H^{q+1}(M;\cI^2) \to


\dotsb


$$


откуда $H^q(M;\cT) \cong H^q(M;\wt\cF_b)$, $q> 0$.


Так как $\wt\cF_b = (\cF_b)^X$, получаем $H^q(M;\wt\cF_b) \cong


H^q(\Sigma,\cF_b)$ (\cite{Bre}, теорема 10.1б, с.\,59), что доказывает


следствие.


\end{pf}








\begin{cor} Для симплектического многообразия с особенностями Мартине


такого, что $\Sigma'=\emptyset$,


имеет место изоморфизм групп когомологий


$$


H^q(M;\gX_h) \cong H^{q+1}_{DR}(M), \quad q>2,


$$


где $\gX_h$~--- пучок гамильтоновых векторных полей на $M$,


$H^k_{DR}(M)$~--- когомологии де Рама многообразия $M$.


\end{cor}





\begin{pf}


Если $\Sigma'=\emptyset$, то одномерное слоение $\cF$ не имеет


особенностей. Тогда для пучка базовых функций $\cF_b$ есть тонкая


резольвента


$$


0 \to \cF_b \to \Omega^0_{\cF} \to \Omega^1_{\cF} \to 0,


$$


где $\Omega^k_{\cF}$~--- пучки форм на $T\cF$ \cite{Vais}.


Следовательно, $H^q(M;\cT) \cong H^q(M;\cF_b) \cong 0$ при $q > 1$.


Теперь наше утверждение вытекает из существования точной


последовательности (\ref{eq7}).


\end{pf}








\begin{cor} \label{cor1}


Если слоение $\cF$ есть расслоение, то $H^q(M;\gX_\Gamma) \cong


H^{q+1}_{DR}(M)$ при $q>0$. В частности, если $H^2_{DR}(M) \cong 0$,


то $\omega$ инфинитезимально жесткая.


\end{cor}





\begin{pf}


Если $\cF$ есть расслоение, то $H^q(M;\cT) \cong H^q(M;\cF_b)\cong 0$


при $q>0$. Требуемое утверждение вытекает из существования точной


последовательности (\ref{eq7}).


\end{pf}








\begin{examnonum}


Пусть $(\wh M,\wh\omega)$~--- симплектическое многообразие, $M$~---


многообразие той же размерности что и $\wh M$, и


$\pi: M \to \wh M$~--- гладкое отображение. Пусть $\omega =


\pi^*\omega$.


Тогда $\omega$~--- замкнутая 2-форма на $M$ и имеет место


коммутативная диаграмма


$$


\begin{CD}


TM @>{\alpha_\omega}>> T^*M


\\


@V{d\pi}VV @AA{d\pi^*}A


\\


T\wh M @>{\alpha_{\Hat \omega}}>> T^*\wh M.


\end{CD}


$$


Из нее следует


\begin{equation}


\label{eq16}


\Ann(\omega) = \ker\alpha_\omega =


d\pi^{-1}(\alpha^{-1}_{\Hat\omega}\ker d\pi^*).


\end{equation}





Пусть теперь $\wh M$ есть $\R^4$ (с координатами $y^\alpha$,


$\alpha=\ov{1,4}$),


наделенное стандартной симплектической структурой


$$


\wh\omega = dy^1\wedge dy^2 + dy^3 \wedge dy^4,


$$


$M = \SS^4 \subset \R^5$~--- стандартная сфера с уравнением


$\sum\limits_{i=1}^5 (x^i)^2 = 1$ и $\pi : \SS^4 \to \R^4$ есть


ограничение проекции $\R^5 \to \R^4$,


$y^\alpha = x^\alpha$, $\alpha=\ov{1,4}$.





Пусть $\Pi$ есть гиперплоскость в $\R^5$ с уравнением $x^5=0$ и $\SS^3_0 =


\SS^4 \cap


\Pi$~--- экватор сферы $\SS^4$.


Ясно, что $d\pi_p : T_p\SS^4 \to T_p\R^4$ есть изоморфизм для любой


$p \in \SS^4 \setminus \SS^3_0$, поэтому $\omega$ невырождена


на $\SS^4 \setminus \SS^3_0$. Далее, вычисление в координатах с


использованием (\ref{eq16}) показывает, что в точках $\SS^3_0$


$\Ann(\omega)$ натянут на векторные поля $\pd_5$ и


$-x^2\pd_1 + x^1\pd_2 - x^4\pd_3 + x^3\pd_4$. Следовательно,


$\Ann(\omega) \pitchfork \SS^3_0$ и одномерное слоение $\cF$,


соответствующее распределению $\Ann(\omega) \cap T\SS^3_0$ на


$\SS^3_0$, есть расслоение Хопфа. Поэтому $\omega$ есть


симплектическая структура с особенностями Мартине типа I, и выполнены


условия следствия \ref{cor1}. Значит, структура $\omega$ является


жесткой.


\end{examnonum}








\begin{thebibliography}{99}





\bibitem{MA} %1


Малахальцев М.А. {\em Инфинитезимальные деформации симплектической


структуры с особенностями} // Изв. вузов. Математика. -- 2003. --


\No\,11. -- C.\,42--50.


\bibitem{Mart1} %2


Martinet J. {\em Sur les singularit\'es des formes diff\'erentielles} //


Ann. Inst. Fourier (Grenoble). -- 1970. -- V.\,20. -- \No\,1. -- P.\,95--178.


\bibitem{Mart2} %3


Martinet J. {\em Singularities of smooth functions and maps} // London


Math. Soc. Lecture Note Ser. --1982. -- \No\,58. -- 256 p.


\bibitem{Rous} %4


Roussarie R. {\em Mod\`eles locaux de champs et de formes} // Ast\'erisque.


-- 1975. -- \No\,30. -- 181 p.


\bibitem{GolTish} %5


Golubitsky M., Tischler D. {\em An example of moduli for singular


symplectic forms} // Invent. Math. -- 1977. -- V.\,38. -- \No\,3. -- P.\,219--225.


\bibitem{SGW} %6


Cannas da Silva A., Guillemin V., Woodwart C. {\em On the unfolding of


folded symplectic structures} // Math. Res. Lett. -- 2000. -- V.\,7.


-- \No\,1. -- P.\,35--53.


\bibitem{S} %7


Cannas da Silva A. {\em Fold-forms for four-folds} // Preprint. -- 2003.


-- 16 p.


\bibitem{D1} %8


Domitrz W. {\em Non-local invariants of Martinet's singular symplectic


structure} //


Banach Center Publ. -- 2002. -- V.\,60. -- P.\,122--143.


\bibitem{Fom} %9


Фоменко А.Т. {\em Симплектическая геометрия. Методы и приложения}.


-- М.: Изд-во МГУ, 1988. -- 414 с.


\bibitem{Bre} %10


Бредон Г. {\em Теория пучков}. -- М.: Наука, 1988. -- 312 c.


\bibitem{Vais}


Vaisman I. {\it $d_f$-cohomologies of Lagrangian


foliations} //


Manatschefte f\"ur Mat. -- 1988 -- V.\,106. -- P.\,221--244.





\end{thebibliography}





\vskip 2cm





\post{\it Казанский государственный}{\it Поступила}


\post{\it университет}{20.09.2004}





\label{end}


\end{document}


